\documentclass[12pt]{article}
\usepackage[utf8]{inputenc}
\usepackage[spanish]{babel}
\usepackage{amsmath, amssymb}
\usepackage{geometry}
\geometry{margin=2.5cm}

\title{Soluciones Modelo Examen Primer Parcial}
\author{}
\date{}

\begin{document}
\maketitle

\section*{Solución Ejercicio 1}

Queremos calcular
\[
\lim_{x\to 0}\frac{\sqrt{1+\sen x}-e^{x/2}}{x^2}.
\]

Al sustituir \(x=0\), obtenemos una indeterminación del tipo
\[
\frac{0}{0},
\]
pues
\[
\sqrt{1+\sen 0}-e^{0/2}=\sqrt{1}-1=0.
\]

Aplicamos la regla de L'Hôpital:
\[
\lim_{x\to 0}\frac{\sqrt{1+\sen x}-e^{x/2}}{x^2}
=
\lim_{x\to 0}
\frac{\dfrac{\cos x}{2\sqrt{1+\sen x}}-\dfrac12 e^{x/2}}{2x}.
\]

Sustituyendo de nuevo \(x=0\), sigue apareciendo una indeterminación:
\[
\frac{\frac12-\frac12}{0}=\frac00.
\]

Aplicamos L'Hôpital por segunda vez:
\[
\lim_{x\to 0}
\frac{\dfrac{\cos x}{2\sqrt{1+\sen x}}-\dfrac12 e^{x/2}}{2x}
=
\lim_{x\to 0}
\frac{
\left(\dfrac{\cos x}{2\sqrt{1+\sen x}}\right)' - \dfrac14 e^{x/2}
}{2}.
\]

Calculamos la derivada del primer término:
\[
\frac{\cos x}{2\sqrt{1+\sen x}}
=
\frac12 \cos x \,(1+\sen x)^{-1/2}.
\]

Entonces
\[
\left(\frac{\cos x}{2\sqrt{1+\sen x}}\right)'
=
\frac12\left[
-\sen x\,(1+\sen x)^{-1/2}
+
\cos x\left(-\frac12\right)(1+\sen x)^{-3/2}\cos x
\right].
\]

Es decir,
\[
\left(\frac{\cos x}{2\sqrt{1+\sen x}}\right)'
=
-\frac{\sen x}{2\sqrt{1+\sen x}}
-\frac{\cos^2 x}{4(1+\sen x)^{3/2}}.
\]

Por tanto,
\[
\lim_{x\to 0}\frac{\sqrt{1+\sen x}-e^{x/2}}{x^2}
=
\lim_{x\to 0}
\frac{
-\dfrac{\sen x}{2\sqrt{1+\sen x}}
-\dfrac{\cos^2 x}{4(1+\sen x)^{3/2}}
-\dfrac14 e^{x/2}
}{2}.
\]

Ahora ya podemos sustituir \(x=0\):
\[
-\frac{\sen 0}{2\sqrt{1+\sen 0}}=0,
\qquad
-\frac{\cos^2 0}{4(1+\sen 0)^{3/2}}=-\frac14,
\qquad
-\frac14 e^{0}=-\frac14.
\]

Luego
\[
\lim_{x\to 0}\frac{\sqrt{1+\sen x}-e^{x/2}}{x^2}
=
\frac{-\frac14-\frac14}{2}
=
\frac{-\frac12}{2}
=
-\frac14.
\]

\[
\boxed{
\lim_{x\to 0}\frac{\sqrt{1+\sen x}-e^{x/2}}{x^2}
=
-\frac14
}
\]

\section*{Solución Ejercicio 2}

Se considera la función
\[
f(x)=
\begin{cases}
\dfrac{e^x-1}{x}, & x<0,\\[8pt]
a, & x=0,\\[8pt]
\dfrac{\sqrt{1+2x}-1}{x}+bx, & x>0,
\end{cases}
\qquad a,b\in\mathbb{R}.
\]

Queremos determinar los valores de \(a\) y \(b\) para que la función sea continua y derivable en \(x=0\).

\subsection*{1. Continuidad en \(x=0\)}

Para que \(f\) sea continua en \(x=0\), debe cumplirse
\[
\lim_{x\to 0^-}f(x)=f(0)=\lim_{x\to 0^+}f(x).
\]

\subsubsection*{Límite por la izquierda}

Si \(x<0\),
\[
f(x)=\frac{e^x-1}{x}.
\]

Usamos el infinitesimo $e^x-1 \sim x$
\[
\lim_{x\to 0}\frac{e^x-1}{x}=1.
\]

Por tanto,
\[
\lim_{x\to 0^-}f(x)=1.
\]

\subsubsection*{Límite por la derecha}

Si \(x>0\),
\[
f(x)=\frac{\sqrt{1+2x}-1}{x}+bx.
\]

Racionalizamos:
\[
\frac{\sqrt{1+2x}-1}{x}
=
\frac{(\sqrt{1+2x}-1)(\sqrt{1+2x}+1)}{x(\sqrt{1+2x}+1)}
=
\frac{2x}{x(\sqrt{1+2x}+1)}
=
\frac{2}{\sqrt{1+2x}+1}.
\]

Así,
\[
f(x)=\frac{2}{\sqrt{1+2x}+1}+bx.
\]

Entonces
\[
\lim_{x\to 0^+}f(x)
=
\frac{2}{1+1}+0
=
1.
\]

Para que la función sea continua en \(x=0\), debe ser
\[
a=1.
\]

\[
\boxed{a=1}
\]

\subsection*{2. Derivabilidad en \(x=0\)}

Suponiendo ya \(a=1\), estudiamos la derivabilidad en \(x=0\).

La función será derivable en \(0\) si existen las derivadas laterales y coinciden:
\[
f'_-(0)=f'_+(0).
\]

\subsubsection*{Derivada lateral izquierda}

Por definición,
\[
f'_-(0)=\lim_{h\to 0^-}\frac{f(h)-f(0)}{h}.
\]

Como para \(h<0\),
\[
f(h)=\frac{e^h-1}{h},
\qquad
f(0)=1,
\]
tenemos
\[
f'_-(0)
=
\lim_{h\to 0^-}\frac{\dfrac{e^h-1}{h}-1}{h}
=
\lim_{h\to 0^-}\frac{e^h-1-h}{h^2}.
\]

Aparece una indeterminación \(\frac00\), así que aplicamos L'Hôpital:
\[
\lim_{h\to 0^-}\frac{e^h-1-h}{h^2}
=
\lim_{h\to 0^-}\frac{e^h-1}{2h}.
\]

Sigue siendo una indeterminación \(\frac00\), así que aplicamos L'Hôpital otra vez:
\[
\lim_{h\to 0^-}\frac{e^h-1}{2h}
=
\lim_{h\to 0^-}\frac{e^h}{2}
=
\frac12.
\]

Por tanto,
\[
f'_-(0)=\frac12.
\]

\subsubsection*{Derivada lateral derecha}

Por definición,
\[
f'_+(0)=\lim_{h\to 0^+}\frac{f(h)-f(0)}{h}.
\]

Como para \(h>0\),
\[
f(h)=\frac{\sqrt{1+2h}-1}{h}+bh,
\qquad
f(0)=1,
\]
resulta
\[
f'_+(0)
=
\lim_{h\to 0^+}\frac{\dfrac{\sqrt{1+2h}-1}{h}+bh-1}{h}.
\]

Reordenamos:
\[
f'_+(0)
=
\lim_{h\to 0^+}
\left(
\frac{\dfrac{\sqrt{1+2h}-1}{h}-1}{h}
+
b
\right).
\]

Así,
\[
f'_+(0)
=
\lim_{h\to 0^+}
\frac{\sqrt{1+2h}-1-h}{h^2}
+b.
\]

Ahora racionalizamos el término
\[
\sqrt{1+2h}-1.
\]
Como
\[
\sqrt{1+2h}-1=\frac{2h}{\sqrt{1+2h}+1},
\]
se tiene
\[
\frac{\sqrt{1+2h}-1}{h}
=
\frac{2}{\sqrt{1+2h}+1}.
\]

Entonces
\[
\frac{\dfrac{\sqrt{1+2h}-1}{h}-1}{h}
=
\frac{\dfrac{2}{\sqrt{1+2h}+1}-1}{h}.
\]

Operamos en el numerador:
\[
\frac{2}{\sqrt{1+2h}+1}-1
=
\frac{2-(\sqrt{1+2h}+1)}{\sqrt{1+2h}+1}
=
\frac{1-\sqrt{1+2h}}{\sqrt{1+2h}+1}.
\]

Por tanto,
\[
\frac{\dfrac{\sqrt{1+2h}-1}{h}-1}{h}
=
\frac{1-\sqrt{1+2h}}{h(\sqrt{1+2h}+1)}.
\]

Sacamos signo menos:
\[
=
-\frac{\sqrt{1+2h}-1}{h(\sqrt{1+2h}+1)}.
\]

Y usando otra vez que
\[
\sqrt{1+2h}-1=\frac{2h}{\sqrt{1+2h}+1},
\]
obtenemos
\[
-\frac{\sqrt{1+2h}-1}{h(\sqrt{1+2h}+1)}
=
-\frac{2h}{h(\sqrt{1+2h}+1)^2}
=
-\frac{2}{(\sqrt{1+2h}+1)^2}.
\]

Así,
\[
f'_+(0)
=
\lim_{h\to 0^+}
\left(
-\frac{2}{(\sqrt{1+2h}+1)^2}+b
\right).
\]

Sustituyendo \(h=0\),
\[
f'_+(0)
=
-\frac{2}{(1+1)^2}+b
=
-\frac{2}{4}+b
=
b-\frac12.
\]

Para que \(f\) sea derivable en \(0\), debe cumplirse
\[
f'_-(0)=f'_+(0),
\]
es decir,
\[
\frac12=b-\frac12.
\]

De aquí,
\[
b=1.
\]

\[
\boxed{b=1}
\]

\subsection*{3. Valor de la derivada}

Con \(a=1\) y \(b=1\), la función es derivable en \(x=0\) y
\[
f'(0)=\frac12.
\]

\[
\boxed{f'(0)=\frac12}
\]

\subsection*{Conclusión}

La función es continua en \(x=0\) si y solo si
\[
\boxed{a=1}.
\]

La función es derivable en \(x=0\) si y solo si
\[
\boxed{a=1,\qquad b=1}.
\]

En ese caso,
\[
\boxed{f'(0)=\frac12}.
\]

\newpage

\section*{Solución Ejercicio 3}

Queremos calcular
\[
\sqrt[5]{34}
\]
con un error menor que \(10^{-5}\), utilizando un desarrollo de Taylor adecuado.

\subsection*{1. Reescritura conveniente}

Buscamos un número próximo cuya raíz quinta sea fácil de calcular. Observamos que
\[
32=2^5,
\]
y además
\[
34=32\left(1+\frac{2}{32}\right)=32\left(1+\frac{1}{16}\right).
\]

Por tanto,
\[
\sqrt[5]{34}
=
\sqrt[5]{32\left(1+\frac{1}{16}\right)}
=
2\sqrt[5]{1+\frac{1}{16}}.
\]

Definimos
\[
f(x)=2(1+x)^{1/5},
\]
de modo que
\[
f\left(\frac{1}{16}\right)=\sqrt[5]{34}.
\]

Desarrollaremos \(f\) en Taylor alrededor de \(x=0\).

\subsection*{2. Derivadas}

Tenemos
\[
f(x)=2(1+x)^{1/5}.
\]

Derivando sucesivamente:
\[
f'(x)=\frac{2}{5}(1+x)^{-4/5},
\]
\[
f''(x)=\frac{2}{5}\left(-\frac{4}{5}\right)(1+x)^{-9/5}
=
-\frac{8}{25}(1+x)^{-9/5},
\]
\[
f^{(3)}(x)
=
-\frac{8}{25}\left(-\frac{9}{5}\right)(1+x)^{-14/5}
=
\frac{72}{125}(1+x)^{-14/5},
\]
\[
f^{(4)}(x)
=
\frac{72}{125}\left(-\frac{14}{5}\right)(1+x)^{-19/5}
=
-\frac{1008}{625}(1+x)^{-19/5},
\]
\[
f^{(5)}(x)
=
-\frac{1008}{625}\left(-\frac{19}{5}\right)(1+x)^{-24/5}
=
\frac{19152}{3125}(1+x)^{-24/5}.
\]

Evaluando en \(x=0\):
\[
f(0)=2,
\qquad
f'(0)=\frac{2}{5},
\qquad
f''(0)=-\frac{8}{25},
\]
\[
f^{(3)}(0)=\frac{72}{125},
\qquad
f^{(4)}(0)=-\frac{1008}{625}.
\]

\subsection*{3. Polinomio de Taylor de grado 4}

El polinomio de Taylor de grado \(4\) centrado en \(0\) es
\[
P_4(x)
=
f(0)+f'(0)x+\frac{f''(0)}{2!}x^2+\frac{f^{(3)}(0)}{3!}x^3+\frac{f^{(4)}(0)}{4!}x^4.
\]

Sustituyendo:
\[
P_4(x)
=
2+\frac{2}{5}x-\frac{8}{25}\frac{x^2}{2}+\frac{72}{125}\frac{x^3}{6}-\frac{1008}{625}\frac{x^4}{24}.
\]

Simplificando:
\[
P_4(x)
=
2+\frac{2}{5}x-\frac{4}{25}x^2+\frac{12}{125}x^3-\frac{42}{625}x^4.
\]

\subsection*{4. Cota del error}

El resto de Lagrange es
\[
R_4(x)=\frac{f^{(5)}(z)}{5!}x^5,
\qquad z\in(0,x).
\]

Para \(x=\frac{1}{16}\),
\[
\left|R_4\left(\frac{1}{16}\right)\right|
=
\left|
\frac{1}{120}\cdot \frac{19152}{3125}(1+z)^{-24/5}\left(\frac{1}{16}\right)^5
\right|.
\]

Como \(z\in\left(0,\frac{1}{16}\right)\), se cumple \(1+z>1\), luego
\[
(1+z)^{-24/5}<1.
\]

Así,
\[
\left|R_4\left(\frac{1}{16}\right)\right|
<
\frac{19152}{3125\cdot 120}\left(\frac{1}{16}\right)^5.
\]

Calculamos:
\[
\frac{19152}{3125\cdot 120}\left(\frac{1}{16}\right)^5
\approx 1.56\cdot 10^{-7}<10^{-5}.
\]

Por tanto, el polinomio \(P_4\) aproxima \(\sqrt[5]{34}\) con error menor que \(10^{-5}\).

\subsection*{5. Aproximación numérica}

Evaluamos en \(x=\frac{1}{16}\):
\[
\sqrt[5]{34}=f\left(\frac{1}{16}\right)\approx P_4\left(\frac{1}{16}\right).
\]

Entonces
\[
P_4\left(\frac{1}{16}\right)
=
2+\frac{2}{5}\cdot\frac{1}{16}
-\frac{4}{25}\left(\frac{1}{16}\right)^2
+\frac{12}{125}\left(\frac{1}{16}\right)^3
-\frac{42}{625}\left(\frac{1}{16}\right)^4.
\]

Calculando:
\[
P_4\left(\frac{1}{16}\right)
=
2+0.025-0.000625+0.0000234375-0.000001640625.
\]

Por tanto,
\[
\sqrt[5]{34}\approx 2.024396796875.
\]

Redondeando,
\[
\boxed{\sqrt[5]{34}\approx 2.02439680}.
\]

\subsection*{Conclusión}

Hemos obtenido la aproximación
\[
\boxed{\sqrt[5]{34}\approx 2.02439680}
\]
con un error menor que
\[
\boxed{10^{-5}}.
\]
\end{document}